% Generated draft: requires mathematical and source review.
\documentclass[11pt]{article}
\usepackage[margin=25mm]{geometry}
\usepackage{fontspec}
\setmainfont{Latin Modern Roman}
\usepackage{amsmath,amssymb,mathtools,unicode-math}
\setmathfont{Latin Modern Math}
\usepackage{enumitem,needspace,xurl,hyperref}
\setlength{\parindent}{0pt}
\setlength{\parskip}{5pt}
\setlength{\emergencystretch}{3em}
\newcommand{\N}{\mathbb{N}}
\newcommand{\Z}{\mathbb{Z}}
\newcommand{\Q}{\mathbb{Q}}
\newcommand{\R}{\mathbb{R}}
\newcommand{\C}{\mathbb{C}}
\newcommand{\F}{\mathbb{F}}
\newcommand{\A}{\mathbb{A}}
\newcommand{\PP}{\mathbb{P}}
\newcommand{\E}{\mathbb{E}}
\newcommand{\eps}{\varepsilon}
\newcommand{\dd}{\,\mathrm{d}}
\begin{document}
\section*{UnsolvedMath ID: 30004391}
\textbf{Source problem code:} catalogue entry
\subsection{Definitions and mathematical statement}
The supplied cleaned formulation uses the fourfold connected sum
\[
N=\mathbb{CP}^{2}\#\mathbb{CP}^{2}\#\mathbb{CP}^{2}\#\mathbb{CP}^{2},\qquad M=(S^{2}\times N)\setminus\{p\}.
\]
Here \(p\in S^{2}\times N\), and the puncture class \(\alpha\in\pi_{5}(M)\) is represented by the boundary sphere of a sufficiently small coordinate ball about \(p\), included in \(M\), with basepoint and orientation choices. This interpretation follows the cleaned statement; the original expression places the connected-sum notation differently and is ambiguous.

To interpret Lipschitz growth, fix a metric on \(M\), the standard round metric on \(S^{5}\), and basepoints. For \(k\in\N\), \(k\geq1\), write
\[
D_{\alpha}(k)=\inf\{\operatorname{Lip}(f):f:S^{5}\to M\text{ is a based Lipschitz map with }[f]=k\alpha\}.
\]
The target metric and the convention concerning based representatives are not specified in the record. In particular, because \(M\) is noncompact, the intended geometry near the puncture must be supplied before this defines the intended quantitative invariant.

\textbf{Puncture-class question.} Determine the asymptotic growth of \(D_{\alpha}(k)\) as \(k\to\infty\). The cleaned record supplies the lower-growth assertion
\[
D_{\alpha}(k)=\omega(k^{1/6}),
\]
which, under the usual little-\(\omega\) convention, means
\[
\lim_{k\to\infty}\frac{D_{\alpha}(k)}{k^{1/6}}=+\infty.
\]
It also describes \(O(k^{1/5})\) as the obvious upper bound. These are supplied claims, not independently verified here. The question requests a sharper lower estimate or an upper estimate of order strictly smaller than \(k^{1/5}\). The latter phrase could mean \(D_{\alpha}(k)=o(k^{1/5})\), or the stronger assertion that there exist \(\beta<1/5\), \(C>0\), and \(k_{0}\) such that
\[
D_{\alpha}(k)\leq Ck^{\beta}\qquad(k\geq k_{0}).
\]
The record does not distinguish these targets. Likewise, the original request for growth strictly faster than \(k^{1/6}\) is already expressed by the cleaned little-\(\omega\) assertion; the intended additional lower-bound improvement is unspecified.

\textbf{Separate formality question.} Let \(Y\) be a simply connected metric space. Interpret the stated linear-nullhomotopy property as follows: for every finite simplicial complex \(X\), equipped with a fixed piecewise Euclidean metric, there exists \(C(X,Y)>0\) such that for every \(L\geq0\) and every nullhomotopic map \(f:X\to Y\) with \(\operatorname{Lip}(f)\leq L\), there exist a point \(y_f\in Y\) and a homotopy \(H:X\times[0,1]\to Y\) satisfying
\[
H(x,0)=f(x),\qquad H(x,1)=y_f,\qquad
\operatorname{Lip}(H)\leq C(X,Y)(L+1).
\]
Here \(X\times[0,1]\) carries a fixed product metric with the interval of length one; the constant is independent of \(L\) and \(f\). Must \(Y\) be formal?

In the usual rational-homotopy meaning, formality means that the commutative differential graded algebra \(A_{\mathrm{PL}}(Y;\Q)\) of rational polynomial forms is connected to \((H^{*}(Y;\Q),0)\) by a zigzag of quasi-isomorphisms. The record does not specify the coefficient field, the topological or finiteness hypotheses on \(Y\), its metric category, or whether the homotopy bound concerns this full product-metric Lipschitz constant rather than another quantitative homotopy convention. These conventions describe an explicit interpretation requiring review, not recovered source hypotheses. The original reference to condition (iii) is not supplied in full; its mention of a formal space is understood in the cleaned record as background to the converse question, rather than an assumption of formality in that question.
\subsection{Short English statement}
For the puncture class in the supplied punctured six-dimensional manifold, determine the least Lipschitz constants of representatives of its integer multiples more sharply than the supplied bounds. Separately, ask whether uniformly linear Lipschitz nullhomotopies for maps from every finite simplicial complex force rational formality. Metric conventions, the intended improvement of the bounds, and the precise category of target spaces remain unresolved.
\subsection{Sources}
\begin{itemize}
\item Local UnsolvedMath registry, \texttt{data/unsolvedmath/problems.json}, with \texttt{catalog/problems.json} as the supplementary registry: supplied permanent ID \(30004391\), source problem code \texttt{catalogue entry}, title \emph{Lipschitz Growth of Puncture Classes and Formality}. The supplied record labels the problem open and its statement \texttt{corrected\_verified}; these are source labels, not independent mathematical verification. No proposer is named; the supplied proposed year is 2020.
\item Supplied report citation: \emph{Manifolds and Groups (2021)}, Oberwolfach Reports, pp.\ 417--463, DOI \texttt{10.4171/owr/2020/7}. \url{https://doi.org/10.4171/owr/2020/7}. The record attributes the collection to workshop open-problem discussions and states that the official report was checked when restoring the formulation. That verification claim is supplied provenance; no source visit is claimed here. The citation year 2021 and the 2020 designation in the DOI are retained as supplied.
\item Supplied literature source: Berdnikov--Manin, \emph{Scalable spaces}. \url{https://arxiv.org/abs/1912.00590}.
\item Supplied literature assessment dated exactly \texttt{2026-08-22}: the puncture-class problem is labeled open, with only the stated \(\omega(k^{1/6})\) lower bound known in the checked literature and the sharper asymptotic distortion unresolved. This is supplied dated curation derived from the Oberwolfach material, not independently verified. The source explicitly requests independent verification of time-sensitive status and cited claims.
\end{itemize}
\end{document}
